\chapter{Extensión y aproximación de funciones continuas}\label{cap:nuevo13}\label{chap:12}\label{chap:28}

Una función definida en una parte de un espacio puede contener información que queremos conservar al ampliar el dominio. La continuidad uniforme permite extenderla a una clausura cuando el espacio de llegada es completo. Tietze resuelve otro problema: extender una función real continua desde un cerrado a todo un espacio métrico, sin exigir continuidad uniforme. La construcción combina separación de cerrados con una serie de correcciones de tamaño decreciente.

Una vez estudiada la extensión, buscaremos aproximaciones uniformes mediante funciones de una forma prescrita. Stone--Weierstrass conecta esa pregunta con las operaciones de un álgebra y la capacidad de distinguir puntos. Los resultados de convergencia uniforme del capítulo~\ref{cap:nuevo11} permiten desarrollar ambas construcciones con control explícito del error.

\section{Separar cerrados mediante funciones}
\label{lectura:12}
\begin{teorema}[Lema de Urysohn para espacios métricos]\label{ch12:urysohn}
Si $A,B$ son cerrados disjuntos de un espacio métrico $X$, existe una función continua $u:X\to[0,1]$ tal que $u|_A=0$ y $u|_B=1$.
\end{teorema}
\begin{proof}
Si $A=\varnothing$, sirve $u=1$; si $B=\varnothing$ y $A\ne\varnothing$, sirve $u=0$. Cuando ambos son no vacíos, definimos
\[
 u(x)=\frac{\dist(x,A)}{\dist(x,A)+\dist(x,B)}.
\]
Las funciones de distancia son continuas, incluso $1$-Lipschitz. El denominador no se anula: si ambas distancias fueran cero, el carácter cerrado de los conjuntos obligaría a $x\in A\cap B$. Por tanto el cociente es continuo y está entre cero y uno. En $A$, el numerador vale cero; en $B$, la distancia a $B$ es cero y la de $A$ es positiva, por lo que el cociente vale uno. Esto también permite separar los cerrados por los abiertos $u^{-1}([0,1/3))$ y $u^{-1}((2/3,1])$, entendiendo los intervalos como abiertos relativos de $[0,1]$.
\end{proof}
\index{Urysohn, lema de}\index{separacion de cerrados@separación de cerrados}
Los cerrados pueden tener distancia mutua cero; no se exige una separación uniforme. Lo que se necesita es que, para cada punto individual, no se anulen simultáneamente las dos distancias.

\begin{lema}[Una corrección que reduce el error]\label{ch12:correccion}
Sea $V$ cerrado en $X$, y sea $g:V\to\R$ continua con $|g|\leq\alpha$, donde $\alpha>0$. Existe una función continua $h:X\to\R$ tal que
\[
 |h(x)|\leq\alpha/3\quad(x\in X),\qquad
 |g(x)-h(x)|\leq2\alpha/3\quad(x\in V).
\]
\end{lema}
\begin{proof}
Los conjuntos
\[
 A=\{x\in V:g(x)\leq-\alpha/3\},\qquad
 B=\{x\in V:g(x)\geq\alpha/3\}
\]
son cerrados en $V$ por continuidad, y por tanto cerrados en $X$ porque $V$ lo es. Son disjuntos. Sea $u$ la función del lema de Urysohn con $u=0$ en $A$ y $u=1$ en $B$, y pongamos $h=(\alpha/3)(2u-1)$. Entonces $|h|\leq\alpha/3$, $h=-\alpha/3$ en $A$ y $h=\alpha/3$ en $B$.

Si $x\in A$, los valores de $g(x)$ están entre $-\alpha$ y $-\alpha/3$, de modo que $|g(x)+\alpha/3|\leq2\alpha/3$. En $B$, $g(x)$ está entre $\alpha/3$ y $\alpha$, y $|g(x)-\alpha/3|\leq2\alpha/3$. En $V\setminus(A\cup B)$ se tiene $|g(x)|<\alpha/3$ y $|h(x)|\leq\alpha/3$, así que la desigualdad triangular da el mismo resultado. Estos tres casos cubren todo $V$.
\end{proof}

Para sumar las correcciones usaremos los resultados del capítulo~\ref{cap:nuevo11}. El criterio M de Weierstrass, teorema~\ref{ch12:weierstrass}, garantiza convergencia uniforme si las cotas de los términos forman una serie numérica convergente, y controla el error por la cola de esa serie. El teorema~\ref{ch12:limite-continuo} asegura después que el límite de las sumas parciales continuas es continuo. En ambos pasos es esencial que el índice que controla el error no dependa del punto. Estos resultados ya están demostrados, también sin exigir acotación global en el teorema de continuidad.

\section{Demostración completa del teorema de extensión}
\begin{teorema}[Tietze]\label{ch12:tietze}
Sea $V$ un cerrado de un espacio métrico $X$. Toda función continua $g:V\to\R$ tiene una extensión continua $f:X\to\R$. Si $g(V)\subset[a,b]$, puede elegirse una extensión con valores en el mismo intervalo.
\end{teorema}
\index{Tietze, teorema de}\index{extension de funciones@extensión de funciones}
\begin{proof}
Si $V=\varnothing$, una función constante resuelve el problema; para la afirmación del intervalo elegimos su valor en $[a,b]$. Si $X=V$, sirve $g$. Tratamos el caso restante.

\textbf{Caso acotado y simétrico.} Supongamos $|g|\leq\alpha$ para algún $\alpha>0$. Aplicamos el lema~\ref{ch12:correccion} a $g$ y obtenemos $h_1$ con $|h_1|\leq\alpha/3$ en $X$ y error a lo sumo $2\alpha/3$ en $V$. El residuo $r_1=g-h_1|_V$ sigue siendo continuo. Aplicamos el mismo lema al residuo con cota $\alpha_1=2\alpha/3$ y obtenemos $h_2$ con
\[
 |h_2|\leq\frac\alpha3\frac23,
 \qquad |g-h_1-h_2|\leq\alpha(2/3)^2\text{ en }V.
\]
Inductivamente, si $h_1,\ldots,h_{n-1}$ ya están construidas, aplicamos el lema al residuo continuo $g-\sum_{k=1}^{n-1}h_k|_V$, cuya cota es $\alpha(2/3)^{n-1}$. Esto produce $h_n$ continua en $X$ con
\[
 |h_n(x)|\leq\frac\alpha3(2/3)^{n-1}\quad(x\in X),
 \qquad
 \left|g(x)-\sum_{k=1}^nh_k(x)\right|\leq\alpha(2/3)^n\quad(x\in V).
\]
La suma del error es \emph{finita} en cada etapa; ése es el residuo al que se aplica el siguiente paso. Como
\[
 \sum_{n=1}^\infty\frac\alpha3(2/3)^{n-1}=\alpha,
\]
el criterio M, teorema~\ref{ch12:weierstrass}, asegura que $f=\sum_nh_n$ converge uniformemente en $X$. Sus sumas parciales son continuas y el teorema~\ref{ch12:limite-continuo} muestra que $f$ es continua. La segunda estimación, al dejar $n\to\infty$, da $f|_V=g$. La primera, sumada, da $|f|\leq\alpha$ en todo $X$.

\textbf{Valores en un intervalo.} Si $a=b$, $g$ es constante y la misma constante la extiende. Si $a<b$, aplicamos el caso anterior a
\[
 g_0=g-\frac{a+b}{2},\qquad |g_0|\leq\frac{b-a}{2}.
\]
Su extensión $f_0$ con la misma cota produce $f=f_0+(a+b)/2$, cuyos valores están en $[a,b]$ y que coincide con $g$ en $V$.

\textbf{Caso no acotado.} La función $h=\arctan g$ es continua en $V$ y toma valores estrictamente entre $-\pi/2$ y $\pi/2$. Por el caso acotado tiene una extensión continua $H:X\to[-\pi/2,\pi/2]$. No podemos aplicar directamente $\tan H$: fuera de $V$, la extensión podría alcanzar los extremos, donde la tangente no está definida. Para evitarlo, definimos
\[
 w(x)=\frac1{1+\dist(x,V)},\qquad \widetilde H(x)=w(x)H(x).
\]
Como $V$ es cerrado y no vacío, $w$ es continua, vale $1$ en $V$ y satisface $0<w(x)<1$ fuera de $V$. En $V$, $|H|<\pi/2$ porque $H=h$; fuera de $V$, el factor $w<1$ garantiza $|\widetilde H|<\pi/2$. Por ello
\[
 f(x)=\tan(\widetilde H(x))
\]
es una función real continua en todo $X$, y en $V$ cumple $f=\tan(\arctan g)=g$. Esto completa el caso que requiere especial cuidado con los extremos del intervalo.
\end{proof}

\begin{ejemplo}[Por qué se necesita un cerrado]
La función $g:(0,\infty)\to\R$, $g(x)=\cos(1/x)$, es continua y está acotada por uno. Sin embargo, para
\[
 x_n=1/(2\pi n),\qquad y_n=1/((2n+1)\pi)
\]
ambas sucesiones tienden a cero y $g(x_n)=1$, $g(y_n)=-1$. Una extensión continua a $\R$ tendría que tener en cero dos límites distintos. La condición de que el dominio original sea cerrado no puede simplemente omitirse.
\end{ejemplo}

\section{Ejercicios sobre separación, soporte y particiones de la unidad}
\begin{ejercicio}[Lectura 12, ejercicio 1]\label{ej:lectura-12-1}
Sin usar el lema de Urysohn, separe dos cerrados disjuntos $A,B\subset X$ por abiertos disjuntos que los contengan.
\end{ejercicio}


\begin{ejercicio}[Lectura 12, ejercicio 2]\label{ej:lectura-12-2}\label{ch12:soporte}
Defina el soporte de una función por
\[
 \supp(f)=\cl{\{x\in X:f(x)\ne0\}}.
\]
Sean $V$ cerrado y $U$ abierto con $V\subset U$.
\begin{enumerate}[label=(\roman*)]
\item Construya $\varphi:X\to[0,1]$ continua, igual a $1$ en $V$, y con $\supp(\varphi)\subset U$.
\item Examine la afirmación de que toda $g:V\to[a,b]$ continua tiene una extensión continua $h:X\to[a,b]$ con soporte en $U$, y resuelva la versión válida.
\end{enumerate}
\end{ejercicio}


\begin{ejercicio}[Lectura 12, ejercicio 3]\label{ej:lectura-12-3}
Sean $U_1,\ldots,U_n$ abiertos que cubren $X$. Construya funciones continuas $\varphi_i:X\to[0,1]$ con $\supp(\varphi_i)\subset U_i$ y $\sum_{i=1}^n\varphi_i=1$.
\end{ejercicio}

\index{particion de la unidad@partición de la unidad}\index{soporte de una funcion@soporte de una función}

Las funciones de separación permiten fijar valores sobre cerrados; las series uniformemente convergentes convierten aproximaciones en una extensión exacta; y las funciones de soporte controlado permiten localizar esa extensión. Una partición de la unidad combina estas ideas para repartir una construcción global entre distintos abiertos sin perder continuidad. Pasaremos ahora de conservar valores en una parte del dominio a aproximar una función en todo un compacto mediante una familia algebraicamente sencilla.

\index{Stone--Weierstrass, teorema de}\index{algebra@álgebra!de funciones}

\section{Cuándo bastan las funciones sencillas}
\label{lectura:28}
El teorema de aproximación de Weierstrass afirma que una función continua en un intervalo compacto puede aproximarse uniformemente por polinomios. Stone--Weierstrass identifica qué propiedades de los polinomios hacen posible esa aproximación: pueden sumarse, multiplicarse y escalarse; distinguen puntos distintos; y no se anulan todos en un mismo punto.

Estudiaremos primero funciones reales. La prueba transformará esas propiedades algebraicas en la posibilidad de construir máximos y mínimos aproximados. Dos aplicaciones de compacidad permitirán ensamblar aproximaciones locales hasta obtener una aproximación en todo el dominio. Para el caso complejo aparecerá una condición adicional: estabilidad bajo conjugación.

Usaremos la norma uniforme, la continuidad de sumas y productos, el teorema de Dini del capítulo~\ref{chap:26} y la compacidad. Construiremos explícitamente los polinomios que aproximan al valor absoluto para evitar utilizar el propio teorema que queremos demostrar.

\section{Álgebras y ejemplos fundamentales}
\begin{definicion}
Un \emph{álgebra asociativa} sobre $\F$ es un espacio vectorial $A$ con una aplicación bilineal $(a,b)\mapsto ab$ de $A\times A$ en $A$ que satisface $(ab)c=a(bc)$. La bilinealidad significa
\[
 (a+b)c=ac+bc,\quad a(b+c)=ab+ac,\quad
 (\lambda a)b=a(\lambda b)=\lambda(ab).
\]
El álgebra es \emph{conmutativa} si $ab=ba$. Es \emph{unitaria} si tiene un elemento $1$ tal que $1a=a1=a$ para todo $a$. Una norma es \emph{submultiplicativa} si $\norm{ab}\le\norm a\norm b$; un álgebra con esa norma es un \emph{álgebra normada}. Si además es completa, es un \emph{álgebra de Banach}.
\end{definicion}
\index{algebra@álgebra!normada}\index{algebra@álgebra!de Banach}

Pedimos que la unidad actúe por ambos lados. Una identidad solamente a la derecha no es, en general, suficiente en un álgebra no conmutativa. Si existen dos unidades bilaterales $e,e'$, entonces $e=ee'=e'$, de modo que la unidad es única.

\begin{ejemplo}[Escalares y matrices]
$\R$ y $\C$ son álgebras conmutativas unitarias sobre sí mismas; $\C$ es también un álgebra sobre $\R$. Las matrices $M_n(\F)$ forman un álgebra con la suma, el producto por escalar y el producto matricial. La asociatividad se puede verificar calculando cada entrada:
\[
 ((AB)C)_{ij}=\sum_k\sum_\ell a_{i\ell}b_{\ell k}c_{kj}
 =\sum_\ell a_{i\ell}\sum_kb_{\ell k}c_{kj}=(A(BC))_{ij}.
\]
La identidad es la matriz $I$. Si $n\ge2$, las matrices elementales satisfacen $E_{12}E_{21}=E_{11}$ y $E_{21}E_{12}=E_{22}$, así que el producto no es conmutativo. Para $n=1$ sí lo es. La norma de operador es submultiplicativa porque $\norm{ABx}\le\norm A\norm B\norm x$.
\end{ejemplo}

\begin{ejemplo}[Cuaternios: un ejemplo asociativo no conmutativo]
\index{cuaternios}
Los cuaternios son las expresiones $q=a+bi+cj+dk$ con coeficientes reales y reglas
\[
 i^2=j^2=k^2=-1,\qquad ij=k,\ jk=i,\ ki=j,
 \qquad ji=-k,\ kj=-i,\ ik=-j.
\]
Estas reglas determinan el producto por bilinealidad. Para comprobar su consistencia y asociatividad sin multiplicar todas las combinaciones de tres elementos, escribimos $z=a+bi$, $w=c+di$ y representamos
\[
 \rho(q)=\begin{pmatrix}z&w\\-\overline w&\overline z\end{pmatrix}.
\]
El conjunto de estas matrices es un subespacio vectorial real y está cerrado bajo producto, pues
\[
 \begin{pmatrix}z&w\\-\overline w&\overline z\end{pmatrix}
 \begin{pmatrix}z'&w'\\-\overline {w'}&\overline {z'}\end{pmatrix}
 =\begin{pmatrix}
 zz'-w\overline {w'}&zw'+w\overline {z'}\\
 -\overline{zw'+w\overline {z'}}&\overline{zz'-w\overline {w'}}
 \end{pmatrix}.
\]
Las matrices de $i,j,k$ satisfacen las reglas indicadas y $\rho$ es inyectiva. La asociatividad del producto de matrices se transmite al producto de cuaternios. El cuaternio conjugado es $\overline q=a-bi-cj-dk$ y la multiplicación da
$q\overline q=\overline q q=(a^2+b^2+c^2+d^2)1$. Así, si $q\ne0$,
\[
 q^{-1}=\frac{\overline q}{a^2+b^2+c^2+d^2}.
\]
Todo elemento no nulo es invertible, pero $ij=-ji$, por lo que esta álgebra de división no es un campo conmutativo. La norma euclidiana es multiplicativa: como $\det\rho(q)=a^2+b^2+c^2+d^2$, la igualdad de determinantes de un producto implica $\norm{qr}_2^2=\norm q_2^2\norm r_2^2$.
\end{ejemplo}

\begin{ejemplo}[Operadores y grupo de invertibles]
Para un espacio normado $V$, los operadores lineales acotados $\mathcal B(V,V)$ forman un álgebra con el producto dado por composición. La asociatividad es la de la composición; las distributividades provienen de la linealidad, y
\[
 \norm{STv}\le\norm S\norm{Tv}\le\norm S\norm T\norm v
 \quad\Longrightarrow\quad\norm{ST}\le\norm S\norm T.
\]
La unidad es $I$. Si $V$ es Banach, el espacio de operadores es completo y constituye un álgebra de Banach.

En cambio,
\[
 \operatorname{GL}(V)=\{T\in\mathcal B(V,V):T\text{ tiene inverso acotado}\}
\]
es un \emph{grupo} bajo composición, no un álgebra cuando $V\ne\{0\}$. En efecto, $I$ y $-I$ son invertibles y su suma es el operador cero, que no es invertible. Ni siquiera es un subespacio vectorial. La cerradura bajo composición se verifica mediante $(ST)^{-1}=T^{-1}S^{-1}$, y ese inverso es acotado por submultiplicatividad. Esto corrige el ejemplo 2(iv) de la lectura original.
\end{ejemplo}

\begin{ejemplo}[Funciones continuas]
Para un espacio métrico $X$, las funciones continuas $C(X,\F)$ forman un álgebra con operaciones puntuales, pues todas las identidades se verifican en cada $x$. Sin embargo, la norma uniforme puede ser infinita si $X$ no es compacto; por ejemplo, $f(x)=x$ en $\R$. El álgebra normada adecuada es $C_b(X,\F)$, donde
\[
 |f(x)g(x)|\le\norm f_\infty\norm g_\infty,
 \qquad\norm{fg}_\infty\le\norm f_\infty\norm g_\infty.
\]
Es un álgebra de Banach por el capítulo~\ref{chap:25}; es conmutativa y tiene por unidad la función constante $1$. Si $X$ es compacto, toda función continua es acotada, y entonces $C(X,\F)=C_b(X,\F)$.
\end{ejemplo}

\begin{definicion}
Una \emph{subálgebra} $A$ de un álgebra $B$ es un subespacio vectorial cerrado bajo multiplicación. No exigimos que contenga la unidad de $B$. Una familia de funciones $A$ \emph{separa puntos} de $X$ si, para $x\ne y$, existe $h\in A$ con $h(x)\ne h(y)$. Decimos que $A$ \emph{no se anula en ningún punto} si para cada $x$ existe $h\in A$ con $h(x)\ne0$.
\end{definicion}
\index{separacion de puntos@separación de puntos}\index{subalgebra@subálgebra}

Los polinomios restringidos a un compacto de $\R$ forman una subálgebra, porque la suma y el producto de polinomios son polinomios. Incluyendo el término de grado cero tenemos la función constante $1$, que garantiza que no hay ceros comunes, y la función $x\mapsto x$ separa puntos. Si solamente se toman potencias de grado positivo y el compacto contiene a cero, todas las funciones de la subálgebra se anulan en ese punto: esa subálgebra no puede ser densa en todo $C(X)$.

\section{Aproximación del valor absoluto sin circularidad}
\begin{lema}\label{ch28:abspolinomios}
Para cada $\alpha\ge0$ existen polinomios reales $p_n$ con $p_n(0)=0$ que convergen uniformemente a $|t|$ en $[-\alpha,\alpha]$.
\end{lema}
\begin{proof}
Si $\alpha=0$, basta $p_n=0$. Construiremos primero polinomios en $[-1,1]$. Definimos
\[
 q_0(t)=0,\qquad q_{n+1}(t)=q_n(t)+\frac{t^2-q_n(t)^2}{2}.
\]
Por inducción, cada $q_n$ es un polinomio y $q_n(0)=0$. Fijemos $t\in[-1,1]$ y escribamos $r=|t|$. Probamos inductivamente que $0\le q_n(t)\le r$ y $q_{n+1}(t)\ge q_n(t)$. Si $0\le q_n(t)\le r$, entonces
\[
 q_{n+1}(t)-q_n(t)=\frac{r^2-q_n(t)^2}{2}\ge0,
\]
y
\[
 r-q_{n+1}(t)
 =(r-q_n(t))\left(1-\frac{r+q_n(t)}2\right)\ge0,
\]
porque $0\le r+q_n(t)\le2$. El caso inicial $q_0=0$ cumple las desigualdades.

Para cada $t$, la sucesión numérica $q_n(t)$ es creciente y está acotada por $r$, por lo que converge a algún $L(t)\in[0,r]$. Al pasar al límite en la recurrencia obtenemos
$L(t)=L(t)+(t^2-L(t)^2)/2$, de modo que $L(t)^2=t^2$. Como $L(t)\ge0$, se sigue que $L(t)=|t|$.

Hemos construido una sucesión creciente de funciones continuas en el compacto $[-1,1]$ cuyo límite $|t|$ es continuo. El teorema de Dini~\ref{ch26:dini} demuestra que la convergencia es uniforme. Para $\alpha>0$, ponemos $p_n(t)=\alpha q_n(t/\alpha)$. Entonces $p_n(0)=0$ y
\[
 \sup_{|t|\le\alpha}|p_n(t)-|t||
 =\alpha\sup_{|u|\le1}|q_n(u)-|u||\longrightarrow0.
\]
\end{proof}

La condición $p_n(0)=0$ tiene una finalidad algebraica: permite escribir $p_n(f)$ usando solo potencias positivas de $f$. Así permaneceremos dentro de una subálgebra aunque no contenga la función constante $1$.

\section{Interpolación en dos puntos}
\begin{lema}\label{ch28:interpolacion}
Si una subálgebra $A$ de funciones $S\to\F$ separa puntos y no se anula en ningún punto, entonces para $x\ne y$ y cualesquiera $\alpha,\beta\in\F$ existe $h\in A$ con $h(x)=\alpha$ y $h(y)=\beta$.
\end{lema}
\begin{proof}
Elegimos $g\in A$ con $g(x)\ne g(y)$, $a\in A$ con $a(x)\ne0$ y $b\in A$ con $b(y)\ne0$. Definimos
\[
 u=gb-g(x)b,\qquad v=ga-g(y)a.
\]
Ambas funciones pertenecen a $A$ porque se obtienen por productos, sumas y múltiplos escalares. No estamos suponiendo que las funciones constantes pertenezcan a $A$: $g(x)b$ es un múltiplo escalar de $b$. Calculando sus valores,
\[
 u(x)=0,\quad u(y)=(g(y)-g(x))b(y)\ne0,
\]
\[
 v(y)=0,\quad v(x)=(g(x)-g(y))a(x)\ne0.
\]
Por tanto,
$h=\alpha v/v(x)+\beta u/u(y)$ pertenece a $A$ y tiene los valores prescritos. No se requiere $\alpha\ne\beta$. Para un solo punto $x$ y un valor $\alpha$, basta elegir $a(x)\ne0$ y usar $\alpha a/a(x)$.
\end{proof}

\section{Demostración completa de Stone--Weierstrass real}
\begin{teorema}\label{ch28:stone}
Sea $X$ un espacio métrico compacto. Si $A$ es una subálgebra de $C(X,\R)$ que separa puntos y no se anula en ningún punto, entonces $\cl A=C(X,\R)$ en la norma uniforme.
\end{teorema}
\begin{proof}
Escribamos $B=\cl A$. Dividimos la prueba en pasos que tienen funciones distintas.

\emph{Paso 1: $B$ es una subálgebra cerrada.}
Para $f,g\in B$ existen sucesiones $f_n,g_n\in A$ con $f_n\to f$ y $g_n\to g$ uniformemente. Las estimaciones
\[
 \norm{(f_n+g_n)-(f+g)}_\infty
 \le\norm{f_n-f}_\infty+\norm{g_n-g}_\infty,
 \qquad\norm{\lambda f_n-\lambda f}_\infty=|\lambda|\norm{f_n-f}_\infty
\]
demuestran que $B$ es un subespacio vectorial. Como $(f_n)$ converge en norma, sus normas están acotadas por algún $M$. Así,
\[
 \norm{f_ng_n-fg}_\infty
 \le M\norm{g_n-g}_\infty+\norm g_\infty\norm{f_n-f}_\infty\longrightarrow0.
\]
Cada $f_ng_n$ pertenece a $A$, de modo que $fg\in B$. La cerradura de $B$ es parte de su definición.

\emph{Paso 2: $B$ contiene valores absolutos y máximos y mínimos finitos.}
Sea $f\in B$ y $\alpha=\norm f_\infty$. El lema~\ref{ch28:abspolinomios} proporciona polinomios $p_n$ sin término constante que aproximan uniformemente al valor absoluto en $[-\alpha,\alpha]$. Cada función $p_n(f)$ pertenece a $B$ porque es combinación lineal de potencias positivas de $f$. Además,
\[
 \norm{p_n(f)-|f|}_\infty
 \le\sup_{|t|\le\alpha}|p_n(t)-|t||\longrightarrow0.
\]
La cerradura de $B$ implica $|f|\in B$. Para $f,g\in B$ utilizamos las identidades puntuales
\[
 \max(f,g)=\frac{f+g+|f-g|}{2},\qquad
 \min(f,g)=\frac{f+g-|f-g|}{2}.
\]
Ambas pertenecen a $B$. Aplicando reiteradamente las identidades se obtienen máximos y mínimos de cualquier familia finita.

\emph{Paso 3: una aproximación global por abajo que coincide en un punto.}
Fijemos $f\in C(X,\R)$, una tolerancia $\eta>0$ y un punto $x\in X$. Para cada $y\in X$, el lema~\ref{ch28:interpolacion} proporciona $h_{x,y}\in A$ con
\[
 h_{x,y}(x)=f(x),\qquad h_{x,y}(y)=f(y).
\]
Si $y=x$ basta la versión de interpolación en un solo punto. El conjunto
\[
 U_y=\{z\in X:h_{x,y}(z)>f(z)-\eta\}
\]
es abierto por continuidad y contiene a $y$. Por compacidad, finitos de ellos, $U_{y_1},\ldots,U_{y_r}$, cubren $X$. Definimos
$g_x=\max\{h_{x,y_1},\ldots,h_{x,y_r}\}\in B$. Todos los términos valen $f(x)$ en $x$, por lo que $g_x(x)=f(x)$. Para cualquier $z\in X$, alguno de los abiertos elegidos contiene a $z$ y entonces
\[
 g_x(z)\ge h_{x,y_j}(z)>f(z)-\eta.
\]
La comparación es con $f(z)$, el valor de la función en el punto variable, y no con la constante $f(x)$.

\emph{Paso 4: controlar también el error por arriba.}
Como $g_x(x)=f(x)$, el conjunto
$V_x=\{z:g_x(z)<f(z)+\eta\}$ es abierto y contiene a $x$. Otra aplicación de compacidad produce una subcubierta $V_{x_1},\ldots,V_{x_m}$ de $X$. Ponemos
\[
 h=\min\{g_{x_1},\ldots,g_{x_m}\}\in B.
\]
Cada $g_{x_j}(z)$ es mayor que $f(z)-\eta$, luego también lo es su mínimo. Por otro lado, para cada $z$ existe $j$ con $z\in V_{x_j}$, y por ello
$h(z)\le g_{x_j}(z)<f(z)+\eta$. Hemos probado
\[
 f(z)-\eta<h(z)<f(z)+\eta\quad(z\in X),
 \qquad\norm{h-f}_\infty\le\eta.
\]

Para concluir, dado $\eps>0$ elegimos $\eta=\eps/3$ en la construcción. Como $h\in B=\cl A$, existe $a\in A$ con $\norm{a-h}_\infty<\eps/3$. Entonces
\[
 \norm{a-f}_\infty\le\norm{a-h}_\infty+\norm{h-f}_\infty
 <2\eps/3<\eps.
\]
Así cada función continua puede aproximarse arbitrariamente por elementos de $A$, como se quería.
\end{proof}

La prueba utiliza máximos y mínimos en la \emph{clausura} del álgebra. No afirma que el máximo de dos polinomios sea un polinomio. La última aproximación de $h\in\cl A$ por un elemento de $A$ devuelve el resultado a la familia original.

\begin{ejemplo}[Qué impide la densidad]
Si todas las funciones de $A$ satisfacen $h(x_0)=0$, todo límite uniforme tiene esa propiedad, pues la evaluación es continua. La función constante $1$ no puede estar en la clausura. Si todas satisfacen $h(x_0)=h(y_0)$ para dos puntos distintos, sus límites uniformes también satisfacen esa igualdad. En un espacio métrico, la función continua $z\mapsto d(z,x_0)$ toma valores distintos en esos dos puntos y, si el espacio es compacto, pertenece a $C(X)$. Por ello la falta de separación impide densidad. Estas observaciones explican las dos hipótesis algebraicas del teorema.
\end{ejemplo}

\section{Ejercicios sobre aproximación polinómica y álgebras complejas}
\begin{ejercicio}[Lectura 28, ejercicio 1: aproximación de Weierstrass]\label{ej:lectura-28-1}
Sea $f:[a,b]\to\F$ continua. Demuestre que existe una sucesión de polinomios con coeficientes en $\F$ que converge uniformemente a $f$.
\end{ejercicio}


\begin{ejercicio}[Lectura 28, ejercicio 2: determinación por momentos]\label{ej:lectura-28-2}
Sea $f:[0,1]\to\R$ continua y suponga que $\int_0^1 f(x)x^n\,dx=0$ para todo entero $n\ge0$. Demuestre que $f=0$.
\end{ejercicio}


\begin{ejercicio}[Lectura 28, ejercicio 3: Stone--Weierstrass complejo]\label{ej:lectura-28-3}
Sea $X$ un espacio métrico compacto. Una subálgebra compleja $A\subset C(X,\C)$ es \emph{autoadjunta} si $\overline f\in A$ para toda $f\in A$. Demuestre que, si $A$ es autoadjunta, separa puntos y no se anula en ningún punto, entonces es densa en $C(X,\C)$.
\end{ejercicio}


\begin{ejercicio}[Lectura 28, ejercicio 4: por qué hace falta conjugación]\label{ej:lectura-28-4}
Sea $\mathbb T=\{z\in\C:|z|=1\}$ y sea $A$ el álgebra de restricciones a $\mathbb T$ de los polinomios complejos en $z$:
\[
 p(e^{i\theta})=\sum_{k=0}^N a_ke^{ik\theta}.
\]
Demuestre que $A$ separa puntos y no se anula en ningún punto, pero no es densa en $C(\mathbb T,\C)$.
\end{ejercicio}


\section{Extender, localizar y aproximar}
Tietze construye una extensión exacta mediante correcciones cuyo error decrece geométricamente; las funciones de soporte controlado permiten localizarla. Stone--Weierstrass resuelve un problema complementario: aproximar una función continua por miembros de una subálgebra. Su argumento crea primero interpolaciones exactas en dos puntos y después las ensambla mediante máximos y mínimos. La aproximación polinómica del valor absoluto permite realizar ese ensamblaje dentro de la clausura de una subálgebra. En el caso complejo, la conjugación recupera partes reales e imaginarias y permite usar el teorema real. El ejercicio~\ref{ej:lectura-28-4} muestra que distinguir puntos no basta si se pierde esa operación. Esta relación entre álgebra y norma uniforme será una herramienta central para las álgebras de operadores y para el análisis armónico de los volúmenes siguientes.
